\paragraph{E.15 (Proved). What the enhanced and marked coefficient categories reconstruct.}
\textit{Manuscript context: LM \S1, folios 1--3, scans 6--11;
Part E(c), the explicit STRATA B.6, B.10 and B.12 leads.}
For small infinity-categories $C,D$, put
$\mathcal P(C)=\operatorname{Fun}(C^{\mathrm{op}},\mathcal S)$.
Then
\[
 \mathcal P(C)\simeq\mathcal P(D)
 \quad\Longleftrightarrow\quad
 \operatorname{Idem}(C)\simeq\operatorname{Idem}(D).
                                                        \tag{E.15.1}
\]
For ordinary categories the right side is equivalence of their ordinary
Karoubi completions. If every endomorphism of $D$ is an equivalence,
an equivalence in (E.15.1) forces every endomorphism of $C$ to be an
equivalence as well. In particular an enhanced presheaf category
remembers this obstruction to being a layered exit category.

For a small space $H$, the ordinary category
$\operatorname{Ho}\operatorname{Fun}(H,\mathcal S)$ together with
its specified constant-diagram functor
$\Delta_H:\operatorname{Ho}\mathcal S\to
\operatorname{Ho}\operatorname{Fun}(H,\mathcal S)$ determines $H$
in $\operatorname{Ho}\mathcal S$. Set-valued local systems together
with their evaluation functors and all natural transformations recover
$\Pi_1H$; finite-valued local systems recover the corresponding
profinite invariant. For connected pointed $H$, an enhanced category
of $R$-module local systems together with its specified evaluation
functor recovers the $E_1$-algebra $R[\Omega H]$.

\begin{proof}
An object $A$ of $\mathcal P(C)$ is \emph{completely compact} when
$\operatorname{Map}(A,-)$ preserves every small colimit. Representables
have this property because mapping out of one is evaluation, and
retracts retain it. Conversely express $A$ as the colimit of its
category of representables. Complete compactness gives
\[
 \operatorname{Map}(A,A)\simeq
    \mathop{\operatorname{colim}}_{(y(c)\to A)}
                       \operatorname{Map}(A,y(c)).
\]
The component of the identity comes from a point of one term:
$\pi_0$ preserves colimits of spaces, and a colimit of sets has no
elements outside the images of its terms. Thus $A$ is a retract of
some $y(c)$. This identifies the intrinsic full subcategory of
completely compact objects with $\operatorname{Idem}(C)$ and proves
the forward implication. The reverse follows from the universal
property of idempotent completion, since $\mathcal S$ is idempotent
complete. This is exactly the reduction in Lurie (2009,
Definition~5.1.6.2, Lemma~5.1.6.4 and Proposition~5.1.6.8)
\cite{FA_LurieHTT2009}, and verifies the supplied STRATA~B.6 proof.
If all endomorphisms are equivalences, every idempotent is an identity
in the homotopy category and splits; hence $D$ is already idempotent
complete. The fully faithful inclusion of $C$ into its completion
then reflects the asserted property of endomorphisms.

For the assertion with specified constants, the intrinsic terminal
object satisfies
\[
 \operatorname{Hom}_{\operatorname{Ho}\operatorname{Fun}(H,\mathcal S)}
                 (1,\Delta_H K)=[H,K]
                                                        \tag{E.15.2}
\]
naturally in $K$. Ordinary Yoneda in the homotopy category of spaces
recovers the represented object $H$. For sets write
$G=\Pi_1H$. The covariant representable $G(x,-)$ represents
evaluation at $x$ on $\operatorname{Fun}(G,\mathbf{Set})$.
Yoneda identifies all transformations between evaluations with the
arrows of $G$, with their composition. This is also consistent with
the intrinsic point-category reconstruction A.8. For a connected
groupoid, automorphisms of evaluation on finite actions form the
inverse limit of the finite quotient groups, explaining precisely
the profinite qualification; A.17 supplies the analogous topological
and pro-groupoid distinctions.

Finally evaluation $U$ at the chosen point of connected $H$ is
conservative and preserves colimits in the enhanced $R$-linear
category. Its left adjoint is
\[
 (LM)(h)=R[\operatorname{Map}_H(x,h)]\otimes_R M.
\]
Consequently the monad $UL$ is tensor product with $R[\Omega_xH]$,
and its multiplication remembers loop composition. Monadicity applies
because $U$ preserves geometric realizations; this is Lurie (2017,
Theorem~4.7.3.5) \cite{FA_LurieHA2017}. Thus the specified enhanced
evaluation recovers that algebra. These arguments independently check
the exact marked assertions of STRATA~B.10.
\end{proof}

The scope of STRATA~B.12 remains the bare \emph{ordinary homotopy}
categories of space-valued or linear coefficient systems, with no
enhancement, constants, or evaluation specified. Neither the
completely compact test, (E.15.2), nor the monad argument reconstructs
that missing structure. This report therefore retains its stated
open qualification for those bare categories, without asserting a
counterexample or a literature-wide open-problem claim. It does not
extend that qualification to the ordinary groupoid-presheaf topoi of
A.8, whose intrinsic point groupoids already give the required
level-one reconstruction. This distinguishes the actual statements
of STRATA~B.6, B.10 and B.12 from the shorter summary in the prompt.
